\section{Frontier topology and a corrected state bound}
\label{sec:frontiers}

The scalar Kauffman-bracket filter is sensitive to the number of matching states
at a scan frontier.  The previous implementation documentation bounded this
number by a Catalan number, using planarity of the diagram.  That argument needs
a geometric hypothesis which the existing greedy ordering does not enforce.
We give a validated one-component counterexample, establish the safe bound for
the current representation, and state precisely when the sharper bound is
available.  This correction concerns a complexity justification; it does not
identify an incorrect bracket evaluation.

\subsection{What the dictionary actually represents}

Fix a planar diagram with $n$ crossings, an order of its crossings, and a nonzero
evaluation point $A\in\mathbb{F}_p$.  After processing $i$ crossings, let $B_i$ be the
set of segment labels occurring exactly once among their PD rows.  These are
the edges with one processed endpoint.  Write $w_i=|B_i|$ and
$w=\max_i w_i$.  Every $w_i$ is even.  Resolving the processed crossings gives
disjoint circles and arcs with endpoints in $B_i$.  The circles contribute
powers of $\delta=-A^2-A^{-2}$; the remaining arcs specify a perfect matching
of $B_i$.

More explicitly, before any final normalization, the partial state is
\begin{equation}
 \sum_{\sigma\in\{0,1\}^i}
 A^{i-2|\sigma|}\delta^{\ell(\sigma)}[M(\sigma)],
 \label{eq:partial-bracket}
\end{equation}
where $\ell(\sigma)$ counts closed circles and $M(\sigma)$ records endpoint
connectivity.  The dictionary combines equal matchings and drops zero
coefficients.  Different embedded arcs with the same labeled connectivity need
not be distinguished: subsequent gluing uses that connectivity to determine
which circles close.  There is no single cyclic order of all frontier points
stored or certified by the greedy scan.

\begin{proposition}[Safe matching count]
\label{prop:frontier-matching-count}
Let $s_i$ be the number of nonzero scalar states after $i$ crossings.  Then
\begin{equation}
 s_i\leq\min\{2^i,(w_i-1)!!\},\qquad (-1)!!=1.
 \label{eq:frontier-matching-count}
\end{equation}
The scan attempts at most $2\sum_{i=0}^{n-1}s_i$ state transitions.
\end{proposition}

\begin{proof}
There are $2^i$ choices of partial resolution.  Every resulting state is a
perfect matching on $w_i$ labeled points.  Pair the smallest available label:
there are $w_i-1$ choices for its partner, then $w_i-3$ choices for the next
pair, and so on.  Hence there are $(w_i-1)!!$ possible matchings.  Combining
coefficients and cancellations only decrease the support.  Each stored state
has two outgoing smoothing transitions at the next crossing.
\end{proof}

With polynomial work per matching operation and the usual dictionary-cost
model, Proposition~\ref{prop:frontier-matching-count} gives
\begin{equation}
 T\leq \poly(n,w,\log p)(w-1)!!
   =2^{O(w\log(w+1))}\poly(n,\log p).
 \label{eq:general-frontier-runtime}
\end{equation}
Ordered dictionaries can provide deterministic comparison guarantees with
additional polynomial factors in the width and label bit lengths.  Arithmetic
counts, dictionary assumptions, and a concrete Python runtime should remain
separate claims.  Persistent matching caches also require their own memory
accounting; a small current coefficient dictionary does not bound every cache
retained from earlier stages.

For even $w$, the exact comparison is
\begin{equation}
 (w-1)!!=\frac{w!}{2^{w/2}(w/2)!}
   =\sqrt{2}(w/e)^{w/2}\bigl(1+O(w^{-1})\bigr).
\end{equation}
This is a safe envelope for the representation, rather than a prediction of
the number of states on a particular input.  The additional $2^i$ bound and
measured support sizes can be much smaller.

\subsection{An explicit knot with more than Catalan-many states}

Consider the following six-crossing PD code, with crossings numbered from zero:
\begin{lstlisting}
[(2,3,0,4), (0,5,1,6), (1,7,8,2),
 (3,9,9,8), (5,10,10,7), (4,11,11,6)]
\end{lstlisting}
Every segment label occurs twice.  The associated rotation system has six
vertices, twelve edges, and eight faces, of degrees
$4,4,6,4,3,1,1,1$, so its connected supporting surface has Euler
characteristic two.  Opposite-port traversal follows the single segment circuit
\[
 0,2,7,10,5,6,11,4,3,9,8,1,0.
\]
Thus this is a spherical diagram with exactly one link component.  The input
also passes the repository's full \code{Diagram.from\_pd} validation.

Process crossings $0,1,2$.  The frontier is
$B_3=\{3,4,5,6,7,8\}$.  Direct expansion of
\eqref{eq:partial-bracket} gives the following six nonzero coefficients in
$\mathbb Z[A,A^{-1}]$.
\begin{center}
\begin{tabular}{ll}
\toprule
Endpoint matching & Laurent coefficient \\
\midrule
$\{(3,5),(4,6),(7,8)\}$ & $2A^{-1}$ \\
$\{(3,5),(4,8),(6,7)\}$ & $A$ \\
$\{(3,6),(4,5),(7,8)\}$ & $A^{-3}+A$ \\
$\{(3,6),(4,8),(5,7)\}$ & $A^{-1}$ \\
$\{(3,8),(4,5),(6,7)\}$ & $A^3$ \\
$\{(3,8),(4,6),(5,7)\}$ & $A$ \\
\bottomrule
\end{tabular}
\end{center}
The third Catalan number is $C_3=5$.  For any fixed cyclic order of six points,
there are exactly five disk-noncrossing perfect matchings.  Consequently no
choice of cyclic order makes all six states disk-noncrossing.  Relabeling the
frontier cannot repair the argument.

This example also survives the actual scalar specialization.  At
\[
 p=2^{61}-1=2305843009213693951,
 \qquad A=2177342782468422681,
\]
all six coefficients remain nonzero.  This matters because distinct smoothing
connectivities alone would not exclude cancellation after aggregation or
specialization.  The JSON artifact records the exact Laurent polynomials, their
field residues, and a witness smoothing for each matching.

The problem is observable under existing ordering code, rather than only a
carefully chosen order.  A second, twelve-crossing validated example has
\code{scan\_order(start=0)} equal to its identity order and reaches the same
six-state frontier.  More directly, the pre-existing
\path{fast/examples/conway_sum_2.json} fixture, under the actual default
\code{best\_scan\_order} with twelve starting choices, begins with
\[
 7,9,4,10,5,8,1,3,0,2.
\]
After this tenth crossing its frontier is
$\{0,5,6,12,13,22\}$, again with six nonzero exact and modular states.
This is a statement about a specified fixture and specified deterministic
ordering, not a statistical claim about typical knots.

The portable verifier
\path{tools/verify_frontier_counterexamples.py} computes every prefix smoothing
from scratch using a connectivity graph on segment labels.  It independently
counts closed components, expands the full Laurent coefficients, and then
compares their field specialization with both the reference gluing routine and
the optimized \code{Planar.glue} routine.  The three checks agree.  The complete
reproducible records are in \path{data/frontier_counterexamples.json}.  The
independent construction is valuable here: using only the same incremental
matching engine to both generate and validate the evidence would provide a
weaker check.

\subsection{The geometric hypothesis that restores Catalan counting}

\begin{proposition}[A certified disk frontier]
\label{prop:disk-frontier}
Suppose a processed tangle is contained in an embedded disk whose boundary
meets the diagram transversely exactly at its frontier points, away from
crossings.  Then its nonzero matching support has cardinality at most
\begin{equation}
 C_{w_i/2}=\frac{1}{w_i/2+1}\binom{w_i}{w_i/2}.
\end{equation}
If this condition holds at every stage, the scalar scan has a
$2^{O(w)}\poly(n,\log p)$ implementation.
\end{proposition}

\begin{proof}
The disk boundary supplies one cyclic order shared by every resolution.  Its
resolved arcs are pairwise disjoint in the disk.  Two pairs with alternating
endpoints cannot occur, by separation in a disk, so every endpoint matching is
noncrossing in that cyclic order.  The number of such matchings on $2k$ points
is $C_k$: pairing a distinguished point divides the remaining points into two
even subintervals and gives the Catalan recurrence.  Since $C_k\leq4^k$, the
state bound and the two smoothing transitions yield the stated runtime, with
the same polynomial per-state accounting as above.
\end{proof}

The hypothesis is sufficient, not asserted necessary.  A regular neighborhood
of a processed crossing set can have several boundary components.  Connectedness
of that set, planarity of the whole diagram, and small cardinality of its
frontier do not produce the single disk required by the proof.  Algorithms
based on sphere-cut decompositions explicitly retain nooses which supply the
missing cyclic boundary information.  The distinction between unrestricted
pairing tables and noose-constrained tables is standard in dynamic programming
on embedded graphs~\cite{ruesauthilikos}; the implementation-specific issue here
is applying the latter count to a scan that does not certify the latter
geometry.

A useful next interface would attach a checkable disk or noose certificate to
an ordering or decomposition.  The checker would verify the embedding and the
boundary incidences independently of coefficient propagation.  A successful
certificate would authorize the Catalan resource estimate; an unsuccessful one
would leave the general matching bound available.  This suggests a concrete
separation of responsibilities: geometry supplies the boundary guarantee, and
the algebraic engine exploits it.

\subsection{What these bounds do and do not say about the target}

For the specific target $n^{O(\log n)}=2^{O(\log^2 n)}$, the general upper
bound~\eqref{eq:general-frontier-runtime} suffices if
\begin{equation}
 w\log(w+1)=O(\log^2 n),
 \label{eq:arbitrary-width-target}
\end{equation}
for example if $w=O(\log^2 n/\log\log n)$ asymptotically.  The certified-disk
bound suffices already at $w=O(\log^2 n)$.  These are sufficient parameter
regimes derived from two different upper bounds.  They are not lower bounds
on the width any successful algorithm must achieve.

In particular, a six-state counterexample does not establish
$(w-1)!!$ lower growth on an infinite family of knot diagrams.  It disproves
the stated Catalan argument and no more.  It does not rule out a sharper
unconditional support theorem, compression of many matching states, or a
different algebraic representation with single-exponential dependence on an
appropriate graph parameter.  Proving a factorial lower bound would require
a family of validated diagrams and orders, together with a proof that the
claimed states occur and remain algebraically independent or nonzero under
the relevant evaluation.  None is claimed here.

Nor does planarity by itself establish the small-width promise in
\eqref{eq:arbitrary-width-target}.  A global recognizer would need a theorem
producing controlled interfaces after permitted reductions, or a route around
large-width residual pieces.  Testing more greedy starts can improve an
individual order, but does not supply that theorem.  This distinction should
guide benchmarks: record the full support profile, frontier widths, cache
costs, and any certified boundary topology, rather than inferring complexity
from crossing number alone.

Finally, all the runtime statements in this section concern a scalar Jones
filter.  A differing value is an obstruction to being the unknot; agreement
is not an unknot certificate supplied by this algorithm.  A fast filter is a
useful part of a complete recognition pipeline, but its favorable width regime
does not resolve the runtime of the residual exact decision procedure.
